% This appendix is a manually curated publication view.
% It intentionally omits repository paths, file names, scripts, and commands.

\chapter{实验参数、评价方法与证据边界}

本附录汇总正文中不宜反复展开、但对学术复核仍有必要的实验设定、数学定义与证据边界。内容按场景设计、故障模型、控制参数、优化约束、评价指标和证据等级组织。变量均在首次出现时说明其意义；仅用于定位代码或执行实验的路径、文件名、脚本名和命令不属于论文论证链，故不在本附录保留。

除特别说明外，表中数值是本文仿真与数字孪生实验的名义设定，不代表传感器制造误差、海试环境统计量或装备验收限值。凡尚未由实物数据标定的参数，均只能用于可控扰动下的相对比较。

\section{仿真场景与扰动因子}\label{sec:app-a-1}

场景库采用“单因素分级 + 多因素联合”的设计。单因素场景用于识别某类观测退化对估计与控制的独立影响，综合压力场景用于检查多种不确定性同时出现时的闭环鲁棒性。令 \(p_{\mathrm{dvl}}\) 表示 DVL 观测在给定时段内的丢失概率，\(B_{\mathrm{sat}}\) 表示磁场硬饱和阈值，\(\eta_a\) 与 \(\eta_s\) 分别表示加速度计和声呐噪声的相对放大系数，\(v_c\) 表示海流速度。

\begin{longtable}{@{}>{\raggedright\arraybackslash}p{0.22\linewidth}>{\raggedright\arraybackslash}p{0.20\linewidth}>{\raggedright\arraybackslash}p{0.20\linewidth}>{\raggedright\arraybackslash}p{0.28\linewidth}@{}}
\caption{仿真场景及其学术用途}\label{tab:app-a-scenarios}\\
\toprule
场景 & 主要扰动 & 名义强度 & 用途 \\
\midrule
\endfirsthead
\toprule
场景 & 主要扰动 & 名义强度 & 用途 \\
\midrule
\endhead
\bottomrule
\endfoot
干净基线 & 无主动扰动 & --- & 给出算法对照基线 \\
DVL 轻度丢失 & 速度观测间歇缺失 & \(p_{\mathrm{dvl}}=0.10\) & 检查轻度观测不连续 \\
DVL 中度丢失 & 速度观测间歇缺失 & \(p_{\mathrm{dvl}}=0.30\) & 检查常见遮挡压力 \\
DVL 重度丢失 & 速度观测间歇缺失 & \(p_{\mathrm{dvl}}=0.60\) & 检查置信度调度边界 \\
DVL 极重度丢失 & 速度观测近乎不可用 & \(p_{\mathrm{dvl}}=0.90\) & 检查估计器失效趋势 \\
轻度磁畸变 & 磁场限幅 & \(B_{\mathrm{sat}}=10^{-6}\,\mathrm{T}\) & 检查磁观测降权机制 \\
重度磁畸变 & 磁场限幅 & \(B_{\mathrm{sat}}=10^{-8}\,\mathrm{T}\) & 检查磁观测近失效状态 \\
声呐杂波 & 成像观测噪声增强 & 固定放大 & 检查声学质量退化 \\
综合压力 & 丢包、漂移、噪声与海流联合 & 见表~\ref{tab:app-a-combined-stress} & 检查多扰动叠加鲁棒性 \\
\end{longtable}

综合压力场景同时施加八类扰动。令 \(p_b\) 表示通信帧丢失概率，\(\dot b_a\) 表示加速度计偏置增长率，\(\lambda_z\) 表示深度尖峰的平均发生率，\(A_z\) 表示尖峰幅值。各量的名义值与物理意图列于表~\ref{tab:app-a-combined-stress}。

\begin{longtable}{@{}>{\raggedright\arraybackslash}p{0.24\linewidth}>{\raggedright\arraybackslash}p{0.28\linewidth}>{\raggedright\arraybackslash}p{0.38\linewidth}@{}}
\caption{综合压力场景的扰动参数}\label{tab:app-a-combined-stress}\\
\toprule
扰动源 & 名义参数 & 物理意图 \\
\midrule
\endfirsthead
\toprule
扰动源 & 名义参数 & 物理意图 \\
\midrule
\endhead
\bottomrule
\endfoot
通信帧丢失 & \(p_b=0.05\) & 近似低带宽链路中的随机丢帧 \\
DVL 丢失 & \(p_{\mathrm{dvl}}=0.30\)，持续 \(1\sim3\,\mathrm{s}\) & 近似浑水或遮挡引起的速度观测中断 \\
IMU 慢漂移 & \(\dot b_a=10^{-3}\,\mathrm{m/s^3}\) & 近似无充分温补时的缓慢偏置增长 \\
深度尖峰 & \(\lambda_z=0.05\,\mathrm{Hz}\)，\(A_z=0.5\,\mathrm{m}\) & 近似低频瞬态电磁干扰 \\
磁场饱和 & \(B_{\mathrm{sat}}=10^{-7}\,\mathrm{T}\) & 近似强近场磁干扰导致的量程受限 \\
加速度计噪声 & \(\eta_a=1.5\) & 近似结构振动耦合 \\
声呐噪声 & \(\eta_s=2.0\) & 近似浑浊浅水中的成像退化 \\
海流 & \(v_c=0.3\,\mathrm{m/s}\) & 近似常见近海横流压力 \\
\end{longtable}

上述磁饱和阈值只用于形成可重复的相对压力，不是由实测噪声谱反演得到的传感器参数。因此，本文只比较同一场景定义下不同算法的相对行为，不把该阈值解释为真实海域中的失效概率。

\section{异步采样与故障注入模型}\label{sec:app-a-2}

仿真接口由通信时延、异步采样和故障注入三个抽象环节组成。通信时延写为
\begin{equation}
  \tau_k=\bar{\tau}+\epsilon_{\tau,k},
  \qquad
  \bar{\tau}=200\,\mathrm{ms},
  \qquad
  \epsilon_{\tau,k}\sim\mathcal{U}(-50,50)\,\mathrm{ms},
  \label{eq:app-comm-delay}
\end{equation}
其中 \(\bar{\tau}\) 为基线时延，\(\epsilon_{\tau,k}\) 为有界抖动。各传感器按独立时钟采样：IMU、深度、磁场和 DVL 的名义频率分别为 \(100\)、\(50\)、\(20\) 和 \(6\,\mathrm{Hz}\)。该设定用于复现多速率观测的时间对齐问题，不代表所有实物传感器的固定输出频率。

六类故障模型定义如下。

\begin{longtable}{@{}>{\raggedright\arraybackslash}p{0.20\linewidth}>{\raggedright\arraybackslash}p{0.35\linewidth}>{\raggedright\arraybackslash}p{0.35\linewidth}@{}}
\caption{故障注入模型及其参数}\label{tab:app-a-chaos}\\
\toprule
故障类型 & 数学或逻辑模型 & 名义参数 \\
\midrule
\endfirsthead
\toprule
故障类型 & 数学或逻辑模型 & 名义参数 \\
\midrule
\endhead
\bottomrule
\endfoot
DVL 冻结 & 丢失窗口内保持最近一次有效速度 \(\tilde{\boldsymbol v}_k=\tilde{\boldsymbol v}_{k_0}\) & 触发概率 \(0.30\)，窗口 \(1\sim3\,\mathrm{s}\) \\
IMU 漂移 & \(\tilde{\boldsymbol a}(t)=\boldsymbol a(t)+\boldsymbol b_a(t)\)，\(\dot{\boldsymbol b}_a=\boldsymbol w_b\) & 名义漂移率 \(10^{-3}\,\mathrm{m/s^3}\) \\
深度尖峰 & \(\tilde z_k=z_k+A_z I_k\)，\(I_k\) 服从泊松到达过程 & \(\lambda_z=0.05\,\mathrm{Hz}\)，\(A_z=0.5\,\mathrm{m}\) \\
磁场饱和 & \(\tilde B_i=\operatorname{clip}(B_i,-B_{\mathrm{sat}},B_{\mathrm{sat}})\) & \(B_{\mathrm{sat}}=10^{-7}\,\mathrm{T}\) \\
通信丢包 & 每帧按伯努利变量 \(I_k\sim\operatorname{Bernoulli}(1-p_b)\) 保留或丢弃 & \(p_b=0.05\) \\
数据乱序 & 以给定概率交换相邻帧的到达顺序 & 乱序概率 \(0.01\) \\
\end{longtable}

这里的 \(\operatorname{clip}(\cdot)\) 表示双边限幅，\(\boldsymbol w_b\) 表示驱动偏置变化的随机过程。模型分别覆盖“持续值错误、缓慢累积误差、瞬态异常、量程受限和通信时序错误”，便于区分估计器对不同失效形态的响应。

\section{不确定性感知控制参数化与消融设计}\label{sec:app-a-3}

令 \(c\in[0,1]\) 为由状态协方差映射得到的归一化置信度，\(c=1\) 表示高置信度。定义低置信度激活量
\begin{equation}
  \rho(c)=(1-c)^\alpha ,
  \label{eq:app-confidence-activation}
\end{equation}
其中 \(\alpha>0\) 控制置信度映射的非线性。跟踪权重和控制增量权重分别写为
\begin{equation}
  w_e(c)=w_{e,0}\left[1+s_e\rho(c)\right],
  \label{eq:app-track-weight}
\end{equation}
\begin{equation}
  w_u(c)=w_{u,0}
  \left[
    \gamma_u+(1-\gamma_u)
    \sigma\!\left(k(c-c_0)\right)
  \right],
  \qquad
  \sigma(x)=\frac{1}{1+\exp(-x)}.
  \label{eq:app-control-weight}
\end{equation}
其中 \(w_{e,0}\) 与 \(w_{u,0}\) 为基准权重，\(s_e\) 为低置信度跟踪权重增益，\(\gamma_u\) 为控制代价的最小缩放比例，\(k\) 为 S 型函数陡峭度，\(c_0\) 为置信度转折点。正文式\eqref{eq:ch04-control-scale}的整体缩放因子 \(s_u(c)\) 即本式方括号项，二者关系为 \(w_u(c)=w_{u,0}\,s_u(c)\)。

\begin{longtable}{@{}>{\raggedright\arraybackslash}p{0.20\linewidth}>{\raggedright\arraybackslash}p{0.17\linewidth}>{\raggedright\arraybackslash}p{0.53\linewidth}@{}}
\caption{不确定性感知控制的名义参数}\label{tab:app-a-uampc-parameters}\\
\toprule
参数 & 名义值 & 含义 \\
\midrule
\endfirsthead
\toprule
参数 & 名义值 & 含义 \\
\midrule
\endhead
\bottomrule
\endfoot
\(s_e\) & \(3.0\) & 低置信度下跟踪权重的最大附加增益 \\
\(\gamma_u\) & \(0.3\) & 低置信度下控制增量代价的最小缩放比例 \\
\(\alpha\) & \(1.5\) & 置信度到跟踪增益的幂指数 \\
\(k\) & \(8.0\) & S 型控制代价缩放的陡峭度 \\
\(c_0\) & \(0.6\) & 高、低置信度过渡的中心值 \\
\(N\) & \(20\) & 模型预测控制的预测步数 \\
\(\Delta t\) & \(0.2\,\mathrm{s}\) & 离散控制时间步长 \\
\end{longtable}

消融实验不以工程开关命名，而按被移除的机制定义。

\begin{longtable}{@{}>{\raggedright\arraybackslash}p{0.12\linewidth}>{\raggedright\arraybackslash}p{0.25\linewidth}>{\raggedright\arraybackslash}p{0.30\linewidth}>{\raggedright\arraybackslash}p{0.23\linewidth}@{}}
\caption{不确定性感知控制消融变体}\label{tab:app-a-uampc-ablation}\\
\toprule
变体 & 定义 & 被隔离的机制 & 学术用途 \\
\midrule
\endfirsthead
\toprule
变体 & 定义 & 被隔离的机制 & 学术用途 \\
\midrule
\endhead
\bottomrule
\endfoot
A0 & \(w_e,w_u\) 恒定 & 全部置信度调度 & 固定权重对照 \\
A1 & 使用完整参数化 & 无 & 默认 UA-MPC \\
A2 & 令 \(w_u(c)=w_{u,0}\) & S 型控制代价降权 & 识别可解性与精度主贡献 \\
A3 & 令 \(\alpha=1\) & 非线性置信度映射 & 检查线性映射的边际影响 \\
A4 & 令 \(s_e=0\) & 低置信度跟踪增益 & 分离跟踪权重放大作用 \\
\end{longtable}

第 5 章的机制归因表明，A2 对可解性和 S 弯横向精度的影响最大；A3 只有边际收益。因此本文保留默认 \(\alpha=1.5\)，不依据单次运动学基准修改默认参数。

\section{代价权重与约束}\label{sec:app-a-4}

名义代价函数写为
\begin{equation}
  J=
  \sum_{k=0}^{N-1}
  \left[
    \boldsymbol e_k^\mathsf{T}\boldsymbol Q(c)\boldsymbol e_k
    +\Delta\boldsymbol u_k^\mathsf{T}
     \boldsymbol R(c)\Delta\boldsymbol u_k
  \right]
  +\boldsymbol e_N^\mathsf{T}\boldsymbol Q_f\boldsymbol e_N ,
  \label{eq:app-cost-function}
\end{equation}
其中 \(\boldsymbol e_k\) 为状态跟踪误差，\(\Delta\boldsymbol u_k\) 为控制增量，\(\boldsymbol Q(c)\) 与 \(\boldsymbol R(c)\) 分别为随置信度变化的跟踪和控制权重矩阵。

\begin{longtable}{@{}>{\raggedright\arraybackslash}p{0.27\linewidth}>{\raggedright\arraybackslash}p{0.18\linewidth}>{\raggedright\arraybackslash}p{0.45\linewidth}@{}}
\caption{代价函数的基准权重}\label{tab:app-a-cost-weights}\\
\toprule
代价项 & 基准权重 & 说明 \\
\midrule
\endfirsthead
\toprule
代价项 & 基准权重 & 说明 \\
\midrule
\endhead
\bottomrule
\endfoot
纵向位置误差 \(e_x\) & \(1.0\) & 水平面位置跟踪 \\
横向位置误差 \(e_y\) & \(1.0\) & 电缆中心线跟踪 \\
深度误差 \(e_z\) & \(5.0\) & 直接关联近底安全裕度 \\
航向误差 \(e_\psi\) & \(3.0\) & 约束路径几何跟随 \\
前向速度误差 \(e_u\) & \(0.5\) & 允许适度速度调整 \\
垂向速度误差 \(e_w\) & \(1.0\) & 限制垂向动态偏差 \\
航向指令增量 & \(0.1\) & 抑制航向指令抖动 \\
深度指令增量 & \(0.1\) & 抑制深度指令跳变 \\
推力指令增量 & \(0.05\) & 抑制推进器频繁调节 \\
\end{longtable}

硬约束用于保证物理可行性，速率约束用于限制相邻预测步的变化，带宽约束用于限制指令相对当前状态的偏离。设 \(u_k\) 为前向速度，\(\psi_{\mathrm{cmd},k}\)、\(z_{\mathrm{cmd},k}\) 和 \(T_{\mathrm{cmd},k}\) 分别为航向、深度和推力指令。

\begin{longtable}{@{}>{\raggedright\arraybackslash}p{0.27\linewidth}>{\raggedright\arraybackslash}p{0.38\linewidth}>{\raggedright\arraybackslash}p{0.25\linewidth}@{}}
\caption{模型预测控制约束}\label{tab:app-a-hard-constraints}\\
\toprule
约束 & 数学表达 & 名义范围 \\
\midrule
\endfirsthead
\toprule
约束 & 数学表达 & 名义范围 \\
\midrule
\endhead
\bottomrule
\endfoot
最低航速 & \(u_k\ge u_{\min}\) & \(u_{\min}=0.1\,\mathrm{m/s}\) \\
航向指令 & \(-\pi\le\psi_{\mathrm{cmd},k}\le\pi\) & \([-\pi,\pi]\,\mathrm{rad}\) \\
深度指令 & \(0\le z_{\mathrm{cmd},k}\le z_{\max}\) & \(z_{\max}=50\,\mathrm{m}\) \\
推力指令 & \(0\le T_{\mathrm{cmd},k}\le100\%\) & 禁止反向推力 \\
深度变化率 & \(\lvert\Delta z_{\mathrm{cmd},k}\rvert\le\Delta z_{\max}\) & \(\Delta z_{\max}=0.5\,\mathrm{m/step}\) \\
航向变化率 & \(\lvert\Delta\psi_{\mathrm{cmd},k}\rvert\le\Delta\psi_{\max}\) & \(\Delta\psi_{\max}=8^\circ/\mathrm{step}\) \\
深度带宽 & \(\lvert z_{\mathrm{cmd},k}-z_{\mathrm{current}}\rvert\le z_b\) & \(z_b=3.0\,\mathrm{m}\) \\
航向带宽 & \(\lvert\psi_{\mathrm{cmd},k}-\psi_{\mathrm{current}}\rvert\le\psi_b\) & \(\psi_b=45^\circ\) \\
\end{longtable}

其中 \(\Delta z_{\mathrm{cmd},k}=z_{\mathrm{cmd},k+1}-z_{\mathrm{cmd},k}\)，航向差采用角度环绕后的最短差值。硬约束、速率约束和带宽约束分别对应“物理边界、短时平滑性和中期可跟随性”，不能相互替代。

\section{数值求解、实时性与回退}\label{sec:app-a-5}

优化问题采用内点法求解。每个控制周期只更新当前状态、参考轨迹和置信度，并以上一周期最优控制序列移位后作为初值。名义配置见表~\ref{tab:app-a-solver}。

\begin{longtable}{@{}>{\raggedright\arraybackslash}p{0.25\linewidth}>{\raggedright\arraybackslash}p{0.22\linewidth}>{\raggedright\arraybackslash}p{0.43\linewidth}@{}}
\caption{数值求解的名义配置}\label{tab:app-a-solver}\\
\toprule
项目 & 名义值 & 说明 \\
\midrule
\endfirsthead
\toprule
项目 & 名义值 & 说明 \\
\midrule
\endhead
\bottomrule
\endfoot
收敛容差 & \(10^{-4}\) & 平衡数值精度与在线计算时间 \\
最大迭代次数 & \(100\) & 防止单次求解无限阻塞控制循环 \\
控制周期 & \(50\,\mathrm{ms}\) & 对应 20 Hz ROS 控制更新；预测离散步长 \(\Delta t=0.2\,\mathrm{s}\) \\
预测时域 & \(N=20\) & 对应约 \(4\,\mathrm{s}\) 的前视窗口 \\
优化变量规模 & 约 \(180\) & 多重打靶下约 \(N(n_x+n_u)=20\times(6+3)\)，状态与控制序列联合优化 \\
\end{longtable}

若求解器在当前周期未返回可接受解，则系统显式进入回退逻辑，而不把回退输出记为求解成功。回退顺序为：优先保持最近一次已验证的安全控制输出；若该输出已过期，则退化到几何制导或安全设定值；当状态新鲜度或安全条件不满足时，由更高优先级的行为树与底层保护接管。

三场景、三种子扫描中，最大迭代次数为 100、200 和 400 时的事件加权回退率分别为 0.699、0.685 和 0.594，当前可观测失败占用分别为 88.41、192.21 和 372.20 ms。400 档的诊断帧数较 100 档减少约 26\%，表明阻塞已经降低控制更新频率。因此迭代上限不是本文宣称的性能优化手段，后续应优先改善初值连续性、约束可行性和终端任务定义。

\section{评价指标与数据契约}\label{sec:app-a-6}

水平 RMSE 与控制变化率定义见式~\eqref{eq:ch05-rmse-xy}和式~\eqref{eq:ch05-control-smoothness}；水平 CEP$_{50}$ 为位置误差模长的中位数。NIS 按观测维数 $m$ 分源统计，滤波假设成立时服从 $\chi_m^2$，维数归一化 NIS 的理论期望为 1。回退率为回退控制周期数除以全部控制周期数，失败周期不得记为求解成功。

巡检须分别记录横偏、路由进度、深度及其不确定度、定位质量、目标身份和执行许可。FWHM 同时报告有效发布率、发布结果条件误差及大误差仍发布比例，避免只报被选中剖面的精度；轨迹误差使用声明区间内全部 Track 样本，含回归修正。通信分别报告到达间隔、应用路径 RTT、板内处理时间和故障响应，不能将其相加为严格单向时延。

\section{核心证据与样本量边界}\label{sec:app-a-7}

表~\ref{tab:app-a-data-sources}合并核心证据。正文保留各实验的具体结果，本表只列样本量和外推限制；不同版本、模型层级和时间窗口不合并计数。

\begin{longtable}{@{}>{\raggedright\arraybackslash}p{0.25\linewidth}>{\raggedright\arraybackslash}p{0.25\linewidth}>{\raggedright\arraybackslash}p{0.40\linewidth}@{}}
\caption{核心实验样本量与证据等级}\label{tab:app-a-data-sources}\\
\toprule
实验对象 & 样本量 & 证据边界\\
\midrule
\endfirsthead
\toprule
实验对象 & 样本量 & 证据边界\\
\midrule
\endhead
\bottomrule
\endfoot
ES-EKF 退化回归 & 8 场景$\times$3 种子；35243 次更新 & 连续性与分源一致性，不证明误差随扰动单调变化\\
时序与鲁棒滤波 & 相位对照 30 次；六策略 90 次 & 离线动力学；时间差膨胀不等于延迟状态更新\\
控制主消融 & 六场景$\times$两模式$\times$三种子，共 36 次 & 代理闭环；平滑收益与横向精度分别报告\\
CBF／原生因子闭环 & CBF 18 次；因子矩阵 20/60 s 各 96 次 & 代理控制；终端退出专项 24 次不证明悬停能力\\
FWHM 稳定性 & 七场景$\times$四方法$\times$200 种子 & 修复前冻结基线，条件误差与发布率同时报告\\
场形拒绝留出 & 每场景 200 个独立留出种子 & 两类已知包络；不保证未知多芯、多源场形\\
当前磁巡检 & 两场形各 3 正例、3 无源对照 & 原生 PVS、理想导航、零流；70 m 核心真值误差\\
ROS 横流／多横切 & 历史横流 8 趟；多横切两方向各 3 种子 & 历史配置、合成身份与流速；运行链分开计数\\
数字孪生验收／先验恢复 & 干净 3 次；中重档各 3 次 & 静态先验、规定窗口；恢复须满足产磁几何\\
磁采集与联合反演 & ADC 四采样率各 10 s；777 对齐点 & 子链灵敏度与相对几何；非整机可溯源埋深\\
Jetson 求解器 & 25 W/8 核独立 clean run & warm/cold 与压力档分开；非整栈长期热稳定\\
PC104 通信与安全 & 300 s 通信；18 项同步故障；1800 s 稳态 & 独立运行、零执行器；非水域闭环\\
\end{longtable}

\section{海缆巡检验收判据}\label{sec:app-a-8}

数字孪生验收采用起点健康、有效窗口、逐帧质量和多次运行聚合四级判据。令 \(s_k\) 为路由进度，\(d_{\perp,k}\) 为横向偏移，\(q_k\in\{0,1\}\) 表示埋深估计是否就绪，\(I_k\in\{0,1\}\) 表示第 \(k\) 帧是否进入验收窗口。有效窗口定义为
\begin{equation}
  I_k=
  \mathbb{I}
  \left(
    q_k=1,\;
    s_k\le50\,\mathrm{m},\;
    |d_{\perp,k}|\le2.0\,\mathrm{m}
  \right).
  \label{eq:app-valid-window}
\end{equation}
起点健康检查使用前 \(30\) 个样本，并要求起始路由进度不超过 \(20\,\mathrm{m}\)、起始横向偏移不超过 \(5\,\mathrm{m}\)。该宽松起点阈值用于排除明显错误初始化，不代替正式验收走廊。

令 \(\sigma_{d,k}\) 为埋深估计标准差，\(\sigma_{\max}\) 为规定上限，则不确定度超限比例为
\begin{equation}
  r_\sigma=
  \frac{
    \sum_k I_k\,
    \mathbb{I}(\sigma_{d,k}>\sigma_{\max})
  }{
    \sum_k I_k
  }.
  \label{eq:app-uncertainty-exceed}
\end{equation}
本文要求 \(r_\sigma\le0.05\)。多次运行聚合至少包含 \(N_{\min}=3\) 次独立运行，通过率定义为
\begin{equation}
  r_{\mathrm{pass}}=
  \frac{N_{\mathrm{pass}}}{N_{\mathrm{run}}},
  \qquad
  r_{\mathrm{pass}}\ge0.67.
  \label{eq:app-pass-rate}
\end{equation}

三次数字孪生运行均满足就绪与通过判据，但“初步验收就绪”只表示数据完整性、窗口有效性和工程阈值在数字孪生中一致，不表示已通过真实海缆检测精度验收。

\section{端到端证据层级与恢复条件}\label{sec:app-a-9}

端到端电缆探测由磁场、导航和声学观测共同驱动，每个时刻输出至少包括横向偏移 \(d_\perp\)、路由进度 \(s\)、埋深估计 \(\hat d\)、埋深标准差 \(\sigma_d\)、置信度 \(c\)、磁信噪比和质量状态。具体程序接口、字段名称和消息通道属于软件文档，不影响本文算法定义，故不在此列出。

设先验电缆点为 \(\boldsymbol p_i\)，受控先验畸变写为
\begin{equation}
  \tilde{\boldsymbol p}_i
  =
  \boldsymbol R(\delta\psi)\boldsymbol p_i
  +\boldsymbol t,
  \label{eq:app-prior-distortion}
\end{equation}
其中 \(\delta\psi\) 为平面旋转偏差，\(\boldsymbol t\) 为平移偏差。在线修正估计 \(\hat{\delta\psi}\) 与 \(\hat{\boldsymbol t}\)，并仅在磁观测满足质量门限时接受更新。该设计保证“恢复”来自新观测对先验的校正，而不是直接放宽验收阈值。

算法级扫描识别敏感性；开环回放检验输入误差的传导；轻量闭环验证接口；六自由度闭环检验受动力学约束的恢复。实物子链只对自身测量范围负责。先验投影连续、物理配准收敛和任务完成分别评价，真值只用于离线评分。

恢复实验的名义垂距为 7.5 m、横切摆幅为 0.6 m、验收走廊半宽为 3.4 m；它们是本组适用条件，不是通用海缆参数。第~\ref{sec:ch05-5-5-11}节给出开关修正的最终对照，算法级 $n=1$ 扫描不并入六自由度多次运行。

\section{代理极端场景的因子化定义}\label{sec:app-a-10}

原生仿真后端尚不能完整耦合电缆几何、近底地形、声磁观测和横流。为验证控制闭环在组合压力下能否运行，本文将代理场景表示为四个因子的组合
\begin{equation}
  \mathcal{S}
  =
  \mathcal{G}
  \times
  \mathcal{T}
  \times
  \mathcal{O}
  \times
  \mathcal{C},
  \label{eq:app-proxy-factorization}
\end{equation}
其中 \(\mathcal{G}\) 表示电缆中心线几何，\(\mathcal{T}\) 表示地形坡度与起伏，\(\mathcal{O}\) 表示声磁观测质量，\(\mathcal{C}\) 表示海流速度和方向。

\begingroup
\small
\setlength{\tabcolsep}{3pt}
\renewcommand{\arraystretch}{1.04}
\begin{longtable}{@{}>{\raggedright\arraybackslash}p{0.23\linewidth}>{\raggedright\arraybackslash}p{0.27\linewidth}>{\raggedright\arraybackslash}p{0.40\linewidth}@{}}
\caption{代理极端场景的主导因子}\label{tab:app-a-proxy-scenarios}\\
\toprule
场景 & 主导因子 & 主要考察问题 \\
\midrule
\endfirsthead
\toprule
场景 & 主导因子 & 主要考察问题 \\
\midrule
\endhead
\bottomrule
\endfoot
连续 S 弯 & \(\mathcal{G}\) & 连续曲率变化下的预瞄与横向跟踪 \\
发卡掉头 & \(\mathcal{G}\) & 大航向变化和终端任务定义 \\
坡面穿越 & \(\mathcal{T}\) & 平面制导与近底高度控制耦合 \\
埋设间断 & \(\mathcal{O}\) & 观测中断下的置信度与回退 \\
横流冲击 & \(\mathcal{C}\) & 外部扰动下的横向保持能力 \\
综合极端 & \(\mathcal{G},\mathcal{T},\mathcal{O},\mathcal{C}\) & 多因素同时作用时的闭环可运行性 \\
\end{longtable}
\endgroup

六类场景分别在固定权重和不确定性感知两种模式下完成单种子冒烟测试，共 \(6\times2=12\) 次运行。该结果只能说明代理压力测试链路能够生成有效控制指标；由于每个组合仅 \(n=1\)，不能据此宣称两种控制模式存在统计显著差异，也不能替代原生声磁耦合数字孪生或海试。

\section{近底安全与原生因子闭环的补充矩阵}\label{sec:app-a-11}

本节汇总正文 第~\ref{sec:ch05-5-7-7}节 只给出结论的两张多种子矩阵，供逐条复核。表~\ref{tab:ch05-cbf-terrain-matrix}给出控制屏障函数（CBF）过滤器在低、中、高三档地形与 PID/MPC 两类控制器下的 18 个约 60 s 运行结果；正文已据此写明 1.5 m 安全违规率与海床穿透率均为 0、最小净空下界 2.70--2.90 m、MPC 求解耗时 p95 上界 23.84 ms 等边界结论。

\begin{longtable}[]{@{}
  >{\raggedright\arraybackslash}p{0.13\linewidth}
  >{\centering\arraybackslash}p{0.08\linewidth}
  >{\centering\arraybackslash}p{0.11\linewidth}
  >{\centering\arraybackslash}p{0.11\linewidth}
  >{\centering\arraybackslash}p{0.12\linewidth}
  >{\centering\arraybackslash}p{0.13\linewidth}
  >{\centering\arraybackslash}p{0.10\linewidth}@{}}
\caption{CBF 地形低/中/高 60 s 三种子复核矩阵结果}\label{tab:ch05-cbf-terrain-matrix}\tabularnewline
\toprule\noalign{}
运行 & 成功次数 & 最小净空下界（m） & 平均净空均值（m） & 违规/穿底最大值 & CBF 激活率/最低缩放 & 求解耗时 p95 最大（ms） \\
\midrule\noalign{}
\endfirsthead
\toprule\noalign{}
运行 & 成功次数 & 最小净空下界（m） & 平均净空均值（m） & 违规/穿底最大值 & CBF 激活率/最低缩放 & 求解耗时 p95 最大（ms） \\
\midrule\noalign{}
\endhead
\bottomrule\noalign{}
\endlastfoot
PID·低 & 3/3 & 2.90 & 3.19 & 0 / 0 & 0.669 / 0.678 & N/A \\
PID·中 & 3/3 & 2.90 & 3.14 & 0 / 0 & 0.473 / 0.563 & N/A \\
PID·高 & 3/3 & 2.80 & 3.16 & 0 / 0 & 0.330 / 0.413 & N/A \\
MPC·低 & 3/3 & 2.90 & 3.22 & 0 / 0 & 0.775 / 0.585 & 23.84 \\
MPC·中 & 3/3 & 2.80 & 3.15 & 0 / 0 & 0.544 / 0.500 & 19.62 \\
MPC·高 & 3/3 & 2.70 & 3.21 & 0 / 0 & 0.463 / 0.180 & 23.62 \\
\end{longtable}

表~\ref{tab:ch05-r22-native-factorial}给出原生因子闭环全矩阵与时序审计。该矩阵覆盖曲线几何、地形/埋设、实测磁噪声回放与横流四个因子的 \(2^4\) 个组合，每组 3 个种子，并分别运行固定权重基线与保守不确定性感知控制，共 96 次、每次 20 s；正文已据此写明 96/96 有效、回退与安全违规均为 0、最坏横向 p95 为 1.716 m、最小净空 2.536 m、求解耗时 p95/max 上界 24.02/38.83 ms 且相对三级周期预算未越界等结论。

\begingroup
\small
\setlength{\tabcolsep}{3pt}
\begin{longtable}[]{@{}
  >{\raggedright\arraybackslash}p{0.16\linewidth}
  >{\centering\arraybackslash}p{0.12\linewidth}
  >{\centering\arraybackslash}p{0.13\linewidth}
  >{\centering\arraybackslash}p{0.12\linewidth}
  >{\centering\arraybackslash}p{0.16\linewidth}
  >{\centering\arraybackslash}p{0.10\linewidth}@{}}
\caption{原生因子闭环全矩阵与时序审计}\label{tab:ch05-r22-native-factorial}\tabularnewline
\toprule\noalign{}
运行口径 & 有效运行 & 最坏横向 p95（m） & 最小净空（m） & 求解耗时 p95 / max（ms） & 周期越界 \\
\midrule\noalign{}
\endfirsthead
\toprule\noalign{}
运行口径 & 有效运行 & 最坏横向 p95（m） & 最小净空（m） & 求解耗时 p95 / max（ms） & 周期越界 \\
\midrule\noalign{}
\endhead
\bottomrule\noalign{}
\endlastfoot
\(2^4\)×3 种子×2 模式，20 s & 96/96 & 1.716 & 2.536 & 24.02 / 38.83 & 0 \\
\end{longtable}
\endgroup

\section{证据边界与剩余验证}\label{sec:app-a-12}

表~\ref{tab:ch05-negative-results}集中列出影响结论外推的限制。各项实测或部署验证可独立推进，不能以更多同源仿真代替真实输入。

\begin{longtable}{@{}>{\raggedright\arraybackslash}p{0.23\linewidth}>{\raggedright\arraybackslash}p{0.33\linewidth}>{\raggedright\arraybackslash}p{0.34\linewidth}@{}}
\caption{主要证据边界与后续验证}\label{tab:ch05-negative-results}\\
\toprule
问题域 & 当前成立范围 & 剩余验证\\
\midrule
\endfirsthead
\toprule
问题域 & 当前成立范围 & 剩余验证\\
\midrule
\endhead
\bottomrule
\endfoot
磁计量与场形 & ADC 子链窄带灵敏度、相对几何拟合及已知模型留出拒绝；未知场形不发布深度 & 可溯源电流/距离、九参数标定、真实背景与同步高度的受控横切\\
状态估计 & 分源 NIS 暴露深度与 DVL 的协方差偏差；保持输出不等于定位准确 & 按自由度整定、时钟斜率与延迟更新、真实 DVL 及绝对水平观测\\
先验恢复与目标身份 & 满足产磁几何时恢复；对称磁峰不自行证明目标身份 & 动态先验、邻源隔离、独立声呐身份和全程排他走廊\\
控制与近底安全 & UA 主收益为平滑；CBF 及因子矩阵安全结果受代理模型约束 & 原生全因子、速度/横流边界与终端保持；核验最大耗时及回退\\
巡检机动与时序 & 有限仿真路线可达；撤权零命令仍会漂移；独立 DDS 连续性未验收 & 目标平台同次整任务、原生采样时间、实测导航/流速及停机包络\\
实物通信与执行 & 零执行器 PC104 通信、watchdog 和同步安全状态链成立；RTT 非单向时延 & 真实电压/DVL、ADC--估计--控制全执行量链、共享时钟、热稳定及水域试验\\
\end{longtable}
